\begin{table}[!t]
    \centering
    \caption{Experimental configuration for the effectiveness--detectability study.}
    \label{tab:experimental_setup}
    \setlength{\tabcolsep}{3.2pt}
    \renewcommand{\arraystretch}{1.04}
    \begin{tabularx}{\columnwidth}{>{\raggedright\arraybackslash}p{1.55cm} X}
        \toprule
        \textbf{Setting} & \textbf{Configuration} \\
        \midrule
        Dataset
            & EMBER 2018 and EMBER 2024 Win64; stratified 20\% subsets preserving the benign--malware label distribution. \\
        Features
            & 2,381 dimensions and 17 trigger features for EMBER 2018; 2,568 dimensions and 19 trigger features for EMBER 2024 Win64. \\
        Classifier
            & LightGBM. The classifier, trigger features, and poisoning protocol are fixed within each condition. \\
        Sampling
            & Random and distribution-based malware-similar benign sampling following PMSB. \\
        Poison rates
            & 1\%, 3\%, and 5\%; the compact main tables report the 1\% setting. \\
        Selectors
            & MinPopulation, CountAbsSHAP, CombinedSHAP, FreqBoundedSigned, LowSignedSHAP, CorrCountAbsSHAP, BenignPrototype, and Quantile95. \\
        Sanitizers
            & Isolation Forest, HDBSCAN, and Spectral Signature; poison recall is the primary detectability measure. \\
        Repeats
            & Two runs are aggregated for the current mean $\pm$ standard-deviation results. \\
        Scope
            & Triggers use problem-space-conservative features. The attack is evaluated in feature space; PE binaries are not modified or executed. \\
        \bottomrule
    \end{tabularx}
\end{table}
